\chapter{实验结果与分析}

\section{实验设置}

本文所有模型均在 CPU 环境下运行，不使用 GPU 和深度学习。分类、回归和聚类均基于 scikit-learn 实现，数值处理和结果整理使用 pandas、numpy 和 scipy，可视化使用 matplotlib。监督学习任务采用训练/测试划分和交叉验证调参，随机过程统一设置 \texttt{random\_state=42}。正文中的指标均来自已有 \texttt{results/} 文件。

\section{探索性分析结果}

图 \ref{fig:eda-risk-distribution} 显示 \texttt{high\_risk\_flag} 的类别分布。高风险样本并非多数类，因此分类任务不能只看 Accuracy；若模型倾向预测低风险，可能在准确率较高的同时漏掉大量高风险样本。

\maybefigure{eda_high_risk_flag_distribution.png}{\texttt{high\_risk\_flag} 类别分布}{fig:eda-risk-distribution}

图 \ref{fig:eda-correlation} 展示数值变量相关性。相关性热力图主要用于识别变量之间的线性关联和潜在冗余，不用于推断因果关系。

\maybefigure{eda_numeric_correlation_heatmap.png}{数值变量相关性热力图}{fig:eda-correlation}

图 \ref{fig:eda-boxplots-risk} 比较不同风险组在设备使用、通知、社交媒体、睡眠与活动变量上的分布差异。这些图可以帮助解释模型为什么关注部分行为变量，但仍然只能作为描述性证据。

\maybefigure{eda_boxplots_by_risk.png}{按高风险标记分组的行为变量箱线图}{fig:eda-boxplots-risk}

图 \ref{fig:eda-category-risk} 展示类别变量下的高风险比例。该图用于观察背景类别与风险标记之间的描述性差异，不能解释为类别本身导致风险变化。

\maybefigure{eda_category_risk_rate.png}{类别变量下的高风险比例}{fig:eda-category-risk}

\section{高风险筛查分类实验}

分类实验比较 Logistic Regression、Random Forest 和 Gradient Boosting，并在调参后对 Gradient Boosting 进行阈值选择。表 \ref{tab:classification-threshold-tuning} 是验证集阈值调优结果，表 \ref{tab:classification-tuned-metrics} 是测试集最终指标。最终选择 threshold=0.14，是因为该策略在验证集上满足 Recall 至少 0.60 的约束，并在测试集上取得 Recall=0.6420、F1=0.5355、PR-AUC=0.5084。

\begin{table}[H]
\centering
\caption{验证集阈值调优结果}
\label{tab:classification-threshold-tuning}
\resizebox{\textwidth}{!}{
\begin{tabular}{llrrrrrr}
\toprule
模型 & 阈值策略 & 阈值 & Precision & Recall & F1 & PR-AUC & Balanced Accuracy \\
\midrule
gradient\_boosting & default\_0\_50 & 0.5000 & 0.6429 & 0.2045 & 0.3103 & 0.4805 & 0.5880 \\
gradient\_boosting & max\_f1 & 0.3500 & 0.5752 & 0.4924 & 0.5306 & 0.4805 & 0.7005 \\
gradient\_boosting & recall\_at\_least\_60\_best\_precision & 0.1400 & 0.4175 & 0.6136 & 0.4969 & 0.4805 & 0.6992 \\
gradient\_boosting & recall\_at\_least\_70\_best\_precision & 0.1200 & 0.2636 & 0.7727 & 0.3931 & 0.4805 & 0.6149 \\
\bottomrule
\end{tabular}
}
\end{table}


\input{tables/classification_tuned_metrics}

图 \ref{fig:classification-metrics} 展示不同阈值策略下的分类指标对比。默认阈值 0.50 的 Precision 较高，但 Recall 偏低；threshold=0.14 则牺牲部分 Precision 以提升高风险样本召回，更符合筛查任务对漏判控制的需求。

\maybefigure{classification_tuned_metrics_comparison.png}{分类阈值策略指标对比}{fig:classification-metrics}

图 \ref{fig:classification-confusion} 为最终阈值下的混淆矩阵，测试集中 $TN=566$、$FP=133$、$FN=63$、$TP=113$。该结果说明模型能识别部分高风险样本，但仍存在误报和漏报，因此只能作为初筛工具，不能作为医学诊断工具。

\maybefigure{classification_final_confusion_matrix.png}{最终分类混淆矩阵}{fig:classification-confusion}

PR 曲线和 ROC 曲线分别见图 \ref{fig:classification-pr} 与图 \ref{fig:classification-roc}。在类别不均衡和高风险筛查场景下，PR-AUC 对正类识别能力更敏感；ROC-AUC 则用于补充观察整体排序能力。

\maybefigure{classification_precision_recall_curve.png}{分类 Precision-Recall 曲线}{fig:classification-pr}

\maybefigure{classification_roc_curve.png}{分类 ROC 曲线}{fig:classification-roc}

\section{数字依赖预测与生产力弱预测}

回归任务以 \texttt{digital\_dependence\_score} 为主线，并将 \texttt{productivity\_score} 作为对照。表 \ref{tab:regression-target-comparison} 显示两个目标的最佳模型均为 Gradient Boosting，但预测效果差异明显：数字依赖预测 $R^2=0.9839$、MSE=3.1471、MAE=0.9982；生产力预测 $R^2=-0.0041$，比简单均值预测更弱。

\begin{table}[H]
\centering
\caption{两个回归目标的最佳模型对比}
\label{tab:regression-target-comparison}
\resizebox{\linewidth}{!}{
\begin{tabular}{llrrrrr}
\toprule
目标变量 & 最佳模型 & MAE & MSE & RMSE & $R^2$ & CV 最优 $R^2$ \\
\midrule
\texttt{productivity\_score} & Gradient Boosting & 7.1671 & 85.2031 & 9.2306 & -0.0041 & 0.0119 \\
\texttt{digital\_dependence\_score} & Gradient Boosting & 0.9982 & 3.1471 & 1.7740 & 0.9839 & 0.9804 \\
\bottomrule
\end{tabular}
}
\end{table}


表 \ref{tab:regression-digital-dependence} 给出数字依赖预测中各模型的测试指标。Gradient Boosting 在该目标上表现最好，说明数字依赖得分与设备使用时长、通知数、解锁次数等行为变量之间具有较强的可预测关系。这里仍需强调，可预测关系不等同于因果关系。

\begin{table}[H]
\centering
\caption{\texttt{digital\_dependence\_score} 回归模型测试集指标}
\label{tab:regression-digital-dependence}
\begin{tabular}{lrrrr}
\toprule
模型 & MAE & MSE & RMSE & $R^2$ \\
\midrule
Gradient Boosting & 0.9982 & 3.1471 & 1.7740 & 0.9839 \\
Ridge & 0.7652 & 3.6541 & 1.9116 & 0.9813 \\
Linear Regression & 0.7651 & 3.6538 & 1.9115 & 0.9813 \\
Random Forest & 1.2680 & 4.3875 & 2.0946 & 0.9776 \\
\bottomrule
\end{tabular}
\end{table}


图 \ref{fig:regression-target-comparison} 和图 \ref{fig:regression-dd-observed} 进一步展示两个回归目标的效果差异。数字依赖预测中真实值与预测值接近对角线，说明拟合效果较好。

\maybefigure{regression_target_comparison.png}{两个回归目标预测效果对比}{fig:regression-target-comparison}

\maybefigure{regression_digital_dependence_observed_vs_predicted.png}{数字依赖得分真实值与预测值对比}{fig:regression-dd-observed}

图 \ref{fig:regression-dd-importance} 显示数字依赖预测的置换重要性。重要变量主要集中在 \texttt{device\_hours\_per\_day}、\texttt{notifications\_per\_day} 和 \texttt{phone\_unlocks} 等行为指标，这与数字依赖得分的业务含义一致。

\maybefigure{regression_digital_dependence_permutation_importance.png}{数字依赖预测置换重要性}{fig:regression-dd-importance}

生产力得分结果见表 \ref{tab:regression-productivity}。该目标的 $R^2$ 为负，说明当前可观测数字行为、生活习惯和背景变量不足以稳定预测生产力得分。这个负结果具有方法价值：它提醒我们不能因为某些目标可以被良好预测，就推断同一套特征对所有结果变量都有效。

\begin{table}[H]
\centering
\caption{\texttt{productivity\_score} 回归模型测试集指标}
\label{tab:regression-productivity-metrics}
\begin{tabular}{lrrrrr}
\toprule
模型 & CV 最优 $R^2$ & MAE & MSE & RMSE & $R^2$ \\
\midrule
gradient\_boosting & 0.0119 & 7.1671 & 85.2031 & 9.2306 & -0.0041 \\
random\_forest & 0.0072 & 7.2088 & 85.9208 & 9.2693 & -0.0125 \\
ridge & 0.0043 & 7.1747 & 85.5884 & 9.2514 & -0.0086 \\
linear\_regression & 0.0002 & 7.1999 & 85.8232 & 9.2641 & -0.0114 \\
\bottomrule
\end{tabular}
\end{table}


\section{数字生活方式聚类画像}

聚类实验只使用数字行为与生活习惯数值特征，不使用背景类别和 outcome 字段训练。表 \ref{tab:clustering-model-comparison} 比较了 KMeans、AgglomerativeClustering 和 GaussianMixture 的内部指标。KMeans 在 $k=3$ 时 Silhouette=0.1860，结合可解释性选择为最终画像方案。

\begin{table}[H]
\centering
\caption{聚类算法与簇数对比摘要}
\label{tab:clustering-model-comparison}
\begin{tabular}{lrrrr}
\toprule
算法 & $k$ & Silhouette & Calinski-Harabasz & Davies-Bouldin \\
\midrule
KMeans & 2 & 0.1808 & 752.8023 & 1.9065 \\
KMeans & 3 & 0.1860 & 611.1770 & 1.7959 \\
KMeans & 4 & 0.1406 & 546.4364 & 1.9615 \\
Agglomerative & 3 & 0.1518 & 495.8938 & 1.8768 \\
GaussianMixture & 2 & 0.1478 & 650.8461 & 2.1599 \\
GaussianMixture & 4 & 0.1225 & 408.1725 & 2.2266 \\
\bottomrule
\end{tabular}
\end{table}


图 \ref{fig:clustering-elbow} 和图 \ref{fig:clustering-silhouette} 分别给出 KMeans 肘部曲线和不同算法/簇数的 Silhouette 对比。Silhouette 数值整体不高，说明聚类结构偏弱，不能把三类画像解释为天然清晰分群。

\maybefigure{clustering_kmeans_elbow.png}{KMeans 肘部曲线}{fig:clustering-elbow}

\maybefigure{clustering_silhouette_by_k.png}{不同聚类算法与簇数的 Silhouette 对比}{fig:clustering-silhouette}

PCA 累计解释方差见图 \ref{fig:pca-explained-variance}。前两个主成分累计解释方差为 0.4241，前五个主成分累计解释方差为 0.7895。PCA 在本文中仅用于降维可视化和辅助理解，不作为因果解释。

\maybefigure{pca_explained_variance.png}{PCA 累计解释方差曲线}{fig:pca-explained-variance}

图 \ref{fig:clustering-pca} 展示聚类结果在前两个 PCA 维度上的投影。不同簇之间存在一定分离趋势，但重叠仍然明显，这与较低 Silhouette 指标一致。

\maybefigure{clustering_lifestyle_pca.png}{数字生活方式聚类的 PCA 可视化}{fig:clustering-pca}

表 \ref{tab:clustering-profiles-compact} 和图 \ref{fig:clustering-heatmap} 展示三类探索性画像。Cluster 0 可概括为高社交媒体使用型，Cluster 1 可概括为高设备依赖型，Cluster 2 可概括为低负荷均衡型。这些命名用于描述簇均值特征，不代表固定人群标签。

\begin{table}[H]
\centering
\caption{数字生活方式聚类画像精简表}
\label{tab:clustering-profiles-compact}
\resizebox{\textwidth}{!}{
\begin{tabular}{rrrrrrrrrl}
\toprule
簇 & 样本数 & 设备时长 & 社交媒体分钟 & 睡眠时长 & 睡眠质量 & 高风险比例 & 生产力均值 & 数字依赖均值 & 建议画像名 \\
\midrule
0 & 447 & 6.9533 & 430.2125 & 7.3816 & 2.7899 & 0.1969 & 65.3375 & 34.3296 & 高社交媒体使用型 \\
1 & 1039 & 11.0538 & 131.8787 & 6.1218 & 1.6953 & 0.3725 & 65.9081 & 51.2425 & 高设备依赖型 \\
2 & 2014 & 5.4711 & 113.4275 & 7.8106 & 3.2137 & 0.1142 & 64.9767 & 29.6962 & 低负荷均衡型 \\
\bottomrule
\end{tabular}
}
\end{table}


\maybefigure{clustering_lifestyle_profile_heatmap.png}{聚类画像标准化特征热力图}{fig:clustering-heatmap}

\section{综合比较与结果边界分析}

综合三类任务可以看到，分类模型在 threshold=0.14 下更适合作为偏召回的高风险初筛工具；数字依赖得分可预测性强，适合作为回归主线；生产力得分在当前特征下预测较弱，应作为负结果分析；聚类结果能够形成有解释价值的探索性画像，但边界不清晰。由于数据集为合成数据，且实验基于相关性和预测关系，本文不作医学诊断、不作个体因果推断，也不建议将模型结果直接用于自动化干预。
